\pdfvariable trailerid {[<C3462026090300000000000000000000><C3462026090300000000000000000000>]}
\documentclass[10pt]{article}
\usepackage[margin=0.76in]{geometry}
\usepackage{amsmath,amssymb,amsthm,microtype,hyperref}
\hypersetup{colorlinks=true,linkcolor=blue,urlcolor=blue}
\providecommand{\CRevisionRound}{2}
\newtheorem{theorem}{Theorem}
\newtheorem{lemma}[theorem]{Lemma}
\newtheorem{remark}[theorem]{Remark}
\newcommand{\RR}{\mathbb R}
\newcommand{\norm}[1]{\lVert#1\rVert}
\newcommand{\calL}{\mathcal L}
\setlength{\emergencystretch}{4em}
\title{A Sharp M-Matrix Threshold for a Two-Dimensional\\Oblique Skorokhod Map}
\author{Route-A source-local certificate HCS-C346}
\date{3 September 2026}

\begin{document}
\maketitle
\begin{abstract}
For oblique regulation in the quadrant with the stated two-axis reflection
matrix $R$,
we prove the \emph{fixed point and sharp M-matrix chamber}:
every c\`adl\`ag input has one and only one regulated path precisely when
$\rho\sigma<1$.
\ifnum\CRevisionRound>0
We add \emph{weighted stability, Picard, and time change}: the natural norm
has contraction factor $\sqrt{\rho\sigma}$, explicit solution stability,
monotone geometric iteration, causality and reparameterization covariance.
\fi
\ifnum\CRevisionRound>1
Finally, \emph{sharp wall, executable audit, and Route-A boundary} close
critical nonuniqueness, critical/supercritical nonexistence, triangular and
post-jump faces without upgrading finite evidence to an all-path proof.
\fi
\end{abstract}

\section{Convention and the sharp existence theorem}
Fix $T<\infty$, $\rho,\sigma\geq0$, and
\[
 R=\begin{pmatrix}1&-\rho\\-\sigma&1\end{pmatrix}.
\]
Let $x=(x_1,x_2)\in D([0,T],\RR^2)$ with $x(0)\geq0$ coordinatewise.
A solution consists of c\`adl\`ag $z,y$ such that
\begin{equation}\label{eq:sp}
 z=x+Ry\geq0,\qquad y(0)=0,\qquad y_i\text{ is nondecreasing},
 \qquad \int_{[0,t]}z_i(s)\,dy_i(s)=0.
\end{equation}
The integrand is the \emph{post-jump} state.  Thus a regulator jump is
supported on an axis after the jump, including simultaneous corner jumps.

For scalar c\`adl\`ag $f$, write
\[
 (\calL f)(t)=\sup_{0\leq s\leq t}[-f(s)]_+.
\]

\begin{lemma}[One-sided regulator]\label{lem:scalar}
If $f(0)\geq0$, the unique nondecreasing c\`adl\`ag $r$, with $r(0)=0$,
for which $g=f+r\geq0$ and $\int g\,dr=0$ is $r=\calL f$.  Moreover,
\begin{equation}\label{eq:Llip}
 \norm{\calL f-\calL h}_\infty\leq\norm{f-h}_\infty.
\end{equation}
\end{lemma}

\begin{proof}
The running supremum is the least nondecreasing correction making $f+r$
nonnegative.  It increases only at a new negative running minimum, where the
post-jump corrected state is zero.  Conversely, feasibility forces every
regulator to dominate $\calL f$.  A first strict excess would be an increase
with positive corrected state, contradicting complementarity.  Finally,
positive part and running supremum are one-Lipschitz, proving
\eqref{eq:Llip}.
\end{proof}

\begin{theorem}[Sharp all-input threshold]\label{thm:main}
Problem~\eqref{eq:sp} has exactly one solution for every admissible input and
every finite horizon if and only if
\begin{equation}\label{eq:chamber}
 \rho\sigma<1.
\end{equation}
Inside this chamber, its regulator is the unique solution of
\begin{equation}\label{eq:fixed}
 y_1=\calL(x_1-\rho y_2),\qquad
 y_2=\calL(x_2-\sigma y_1).
\end{equation}
\end{theorem}

\begin{proof}
With the other coordinate held fixed, each line of~\eqref{eq:sp} is the
scalar problem in Lemma~\ref{lem:scalar}; hence~\eqref{eq:sp} and
\eqref{eq:fixed} are equivalent.  Conversely, at time zero the fixed-point
equations and $x(0)\geq0$ imply
$y_1(0)\leq\rho y_2(0)$ and $y_2(0)\leq\sigma y_1(0)$; since
$\rho\sigma<1$, both values vanish.  Thus the fixed point has the required
initial condition.

First suppose $\rho,\sigma>0$, put $w_1=\sqrt\rho$, $w_2=\sqrt\sigma$,
$q=\sqrt{\rho\sigma}$, and define
\[
 \norm{u}_w=\max_{i=1,2}\frac{\norm{u_i}_\infty}{w_i}.
\]
For the right side $\Phi_x$ of~\eqref{eq:fixed}, inequality~\eqref{eq:Llip}
gives
\[
 \norm{\Phi_x(u)-\Phi_x(v)}_w\leq q\norm{u-v}_w.
\]
Thus~\eqref{eq:chamber} makes $\Phi_x$ a contraction on the complete space of
bounded c\`adl\`ag pairs.  Its unique fixed point is c\`adl\`ag and solves the
problem.  The constant $q$ is the best uniform coefficient: each weighted
cross-coordinate estimate can be attained by a constant difference after a
suitable downward input shift.

If $\rho=0$, solve successively
$y_1=\calL x_1$ and $y_2=\calL(x_2-\sigma y_1)$; if $\sigma=0$, reverse the
order.  This also covers $\rho=\sigma=0$ and proves sufficiency.

For necessity, let $\rho\sigma=1$.  Then $R(\rho,1)^\mathsf T=0$, so zero
input has the distinct solutions
\begin{equation}\label{eq:null}
 y=(\rho h,h),\qquad z=0,
\end{equation}
for every nondecreasing c\`adl\`ag $h$ with $h(0)=0$.  More strongly, whenever
$\rho\sigma\geq1$, an input jumping from zero to $(-1,-1)$ cannot be feasible:
the post-jump values would obey
\[
 y_1\geq1+\rho y_2,\qquad y_2\geq1+\sigma y_1,
\]
and hence $(1-\rho\sigma)y_1\geq1+\rho$, an impossibility.  Therefore no
point on or beyond the wall is well posed for every input.
\end{proof}

\ifnum\CRevisionRound>0
\section{Quantitative stability and path structure}
\begin{theorem}[Stability and constructive convergence]\label{thm:stable}
Assume $\rho,\sigma>0$ and $q<1$.  If inputs $x,x'$ have solutions
$(z,y),(z',y')$, then
\begin{equation}\label{eq:stability}
 \norm{y-y'}_w\leq\frac{\norm{x-x'}_w}{1-q},
 \qquad
 \norm{z-z'}_w\leq\frac{2\norm{x-x'}_w}{1-q}.
\end{equation}
Starting with $y^{(0)}=0$ and setting
$y^{(n+1)}=\Phi_x(y^{(n)})$ gives a coordinatewise increasing sequence and
\begin{align}
 \norm{y^{(n+1)}-y^{(n)}}_w&\leq
 q^n\norm{y^{(1)}}_w,\label{eq:successive}\\
 \norm{y-y^{(n)}}_w&\leq
 \frac{q^n}{1-q}\norm{y^{(1)}}_w.\label{eq:tail}
\end{align}
The solution map is causal, preserves continuity, and commutes with continuous
nondecreasing onto time changes that preserve the endpoints.
\end{theorem}

\begin{proof}
For two inputs, the scalar estimate gives
\[
 \norm{\Phi_x(u)-\Phi_{x'}(v)}_w
 \leq\norm{x-x'}_w+q\norm{u-v}_w.
\]
Evaluate it at the two fixed points for the first inequality in
\eqref{eq:stability}.  Each normalized row of $R$ has weighted absolute sum
$1+q$, whence
\[
 \norm{R(y-y')}_w\leq(1+q)\norm{y-y'}_w.
\]
Using $z-z'=x-x'+R(y-y')$ yields the second inequality.  The map $\Phi_x$ is
order preserving and $0\leq\Phi_x(0)$, so Picard iterates increase.
Contraction gives~\eqref{eq:successive}, and summing its geometric tail gives
\eqref{eq:tail}.

Running suprema use only the past, so uniqueness gives causality.  They map
continuous functions to continuous functions; uniform Picard convergence then
preserves continuity.  If
$\lambda:[0,T']\to[0,T]$ is continuous, nondecreasing, onto, and endpoint
preserving, its initial images satisfy
$\lambda([0,t])=[0,\lambda(t)]$.  Consequently
\[
 \calL(f\circ\lambda)=(\calL f)\circ\lambda.
\]
Compose~\eqref{eq:fixed} with $\lambda$ and use uniqueness.  Flat pieces of
$\lambda$ simply pause both state and regulator.
\end{proof}

\begin{remark}
When one weight vanishes, the weighted norm is not defined.  The triangular
formulas in Theorem~\ref{thm:main} are the correct boundary statement; no
singular limiting norm is smuggled into that face.
\end{remark}
\fi

\ifnum\CRevisionRound>1
\section{Boundary closure and executable receipt}
At the critical wall, equation~\eqref{eq:null} proves nonuniqueness even for
the smoothest possible input, whereas the negative jump proves nonexistence.
These are separate failures.  At normal reflection the solution is
$(\calL x_1,\calL x_2)$; a one-sided coupling is triangular.  The post-jump
convention resolves a step input through the linear complementarity problem
for $d=y-y^-$ and the new state $x+R(y^-+d)$.  Strict inequality in
\eqref{eq:chamber}, the off-diagonal signs, and simultaneous cascades are all
therefore essential.

The machine-readable receipt has six rational matrices, 36 original events,
72 pause-expanded events, 27 Picard rows, three critical nonunique regulators,
three no-solution witnesses, and 693 audited scalar leaves.  An independent
implementation performs 886 checks; a separate SymPy lane closes 5,125 exact
identities; two isolated reproductions are byte-identical; and 68 parser,
semantic, stale-hash and repaired-hash attacks are rejected.  These checks
audit conventions.  Theorems~\ref{thm:main}--\ref{thm:stable}, not finite
enumeration, prove the all-input result.

The nearest workspace owners are scalar Moreau play (C332), one-interface
skew Brownian motion (C266), total-variation flow (C279), and dry-friction
capture (C238).  None owns a two-axis oblique reflection map or its sharp
M-matrix threshold.

The strict evaluation is
\[
 (\mathrm{A0\_FAIL},\mathrm{A1\_FAIL},\mathrm{A2\_FAIL},
  \mathrm{A3\_FAIL},\mathrm{A4\_FAIL}).
\]
Free couplings and driven contacts have no rational-prime taxonomy; $\det R$
is not an Euler product or target determinant; and the dissipative regulator
provides no natural target-zero quantization.  Route A is rejected and Route B
is locked under \texttt{NO\_BAD\_EULER\_OR\_ROOT\_NUMBER}.  No target
arithmetic local data, Euler factor, root number, automorphy, divisor or
counting law, functional equation, zero match, or Hilbert--P\'olya operator is
claimed.

\paragraph{AI use.}
A generative language model assisted drafting and code scaffolding.  The
self-contained proof, independent reconstruction, hostile mutations and
deterministic artifacts define the audit record.
\fi

\section*{Source lineage}
This source-local reconstruction makes no literature-priority claim.
\begin{thebibliography}{9}
\bibitem{HR}
J. M. Harrison and M. I. Reiman, ``Reflected Brownian Motion on an
Orthant,'' \emph{Annals of Probability} 9 (1981), 302--308. DOI:
\href{https://doi.org/10.1214/aop/1176994471}{10.1214/aop/1176994471}.
\bibitem{DI}
P. Dupuis and H. Ishii, ``On Lipschitz Continuity of the Solution Mapping to
the Skorokhod Problem, with Applications,'' \emph{Stochastics and Stochastics
Reports} 35 (1991), 31--62. DOI:
\href{https://doi.org/10.1080/17442509108833688}%
{10.1080/17442509108833688}.
\end{thebibliography}
\end{document}
